\begin{abstract}
Text-to-image models produce rectangular images, whereas a usable game-asset
texture must remain coherent after it is cut into a fragmented UV atlas, mapped
onto a narrow three-dimensional surface, viewed from several directions, and
exported under asset-specific constraints. This report presents \skinsmith, a
human-supervised generative Agent that converts a short theme keyword into a
mapped and reproducible weapon texture. The system introduces three checkpoints:
the user confirms an expanded visual theme, selects one textual art direction,
and selects one artwork only after seeing its source image together with left,
right, and top weapon previews.

The technical method separates a tileable-pattern baseline from a weapon-aware
route. The latter derives a canonical coordinate frame from the target mesh,
constructs per-UV-pixel position and normal maps, fits a selected master artwork
to continuous weapon-space canvases, and then bakes the result into the asset's
existing UV atlas. Candidates are evaluated using texture statistics, a
mesh-derived seam metric, and fixed multi-view renders. Selection first enforces
hard constraints and then applies Pareto dominance and a deterministic objective
order. One local refinement round is permitted, but it is retained only when it
improves the locked score by at least 0.01; otherwise the system rolls back.

In a four-pair formal experiment on an AK-47 asset, the weapon-space route reduced
asset-seam error by a mean of 0.2693 and improved the multi-view score by 0.0601,
winning all four paired comparisons on both measures. A separate selection-policy
ablation showed that an equal weighted score chose an infeasible seam width,
whereas constraint-first selection returned a valid operating point. All four
formal refinement attempts were rejected and rolled back, demonstrating bounded
no-regression behavior rather than guaranteed self-improvement. A real
three-checkpoint run generated four mapped artwork alternatives from the keyword
\emph{dragon}; the user-selected candidate was exported as a 2048 by 2048 RGB
texture with asset-seam error 0.00078 and multi-view score 0.8166. Transfer to an
M4A4 adapter reused the core pipeline without new image generation, but did not
satisfy the AK-47 retention threshold, exposing the asset-specific limits of
portability.

The results support the value of composing and evaluating artwork in the target
asset's geometry rather than treating image generation as the final output. They
also show why human selection, explicit feasibility rules, preserved evidence,
and rollback are necessary in a creative Agent whose automatic metrics remain
imperfect.
\end{abstract}
